% --- Archon physics formalization source begin ---
% archon:physics
% archon:covers PhyXMiniProblems/problem_phyx_mini_0424.lean
% archon:source-report reports/phyx_mini/problem_phyx_mini_0424.source.json

\chapter{Physics problem phyx\_mini\_0424}
\label{ch:PhyXMiniProblems_problem_phyx_mini_0424}

\paragraph{Problem source.}
\#\# Physical scenario

Figure shows a Carnot heat engine driving a Carnot refrigerator.

\#\# Question

Determine \$Q\_3\$.

\#\# Answer choices

- A: A: 1000 J

- B: B: 1500 J

- C: C: 2000 J

- D: D: 2500 J

\#\# Figure caption (auxiliary; use the image as primary evidence)

The image represents a thermodynamic diagram involving a Carnot heat engine and a Carnot refrigerator.

- **Sections and Labels:**
  - **Top Section:** 
    - On the left, the temperature is labeled 600 K.
    - On the right, the temperature is labeled 500 K.
    - Between these, an arrow labeled \textbackslash{}(Q\_1 = 1000 \textbackslash{}, \textbackslash{}text\{J\}\textbackslash{}) indicates heat transfer.

  - **Bottom Section:**
    - On the left, the temperature is labeled 300 K.
    - On the right, the temperature is labeled 400 K.

- **Components and Processes:**
  - **Carnot Heat Engine:**
    - Located on the left side.
    - Heat \textbackslash{}(Q\_1\textbackslash{}) enters from the top (600 K) and exits as \textbackslash{}(Q\_2\textbackslash{}) at the bottom (300 K) with downward arrows.
    - Outputs work labeled \textbackslash{}(W\_\{\textbackslash{}text\{out\}\}\textbackslash{}).

  - **Carnot Refrigerator:**
    - Located on the right side.
    - Heat \textbackslash{}(Q\_3\textbackslash{}) enters from the bottom (400 K) and exits as \textbackslash{}(Q\_4\textbackslash{}) at the top (500 K) with upward arrows.
    - Requires work labeled \textbackslash{}(W\_\{\textbackslash{}text\{in\}\}\textbackslash{}).

- **Connections:**
  - The heat engine and refrigerator are connected through work transfers, \textbackslash{}(W\_\{\textbackslash{}text\{out\}\}\textbackslash{}) from the heat engine to \textbackslash{}(W\_\{\textbackslash{}text\{in\}\}\textbackslash{}) for the refrigerator.

- **Arrows:**
  - Purple arrows indicate the direction of heat and work transfers through the systems.

The diagram illustrates the energy exchange and work process between a heat engine and a refrigerator operating between these temperatures.

\#\# Dataset metadata

Laws of Thermodynamics; Multi-Formula Reasoning, Physical Model Grounding Reasoning

\paragraph{Recorded answer/context.}
D: D: 2500 J

\paragraph{Figure/image path.}
/root/proposal\_for\_physic/science-mango/phyx\_mini\_run/phyx\_data/test\_image/424.png

\paragraph{Formalization target.}
create a compiling Lean file with sorry bodies at `PhyXMiniProblems/problem\_phyx\_mini\_0424.lean`. The Lean declarations must preserve the physical quantities, dimensions or dimensional roles, figure labels, governing-law hypotheses, and final relation expressed by this problem.
Use Mathlib/Physlib names found through LeanExplore where available. If a physics API is missing, introduce faithful local abstractions rather than scalar placeholder aliases.

\begin{theorem}[Physics formalization target]
\label{thm:physics:phyx_mini_0424:target}
\lean{PhyXMiniProblems.ProblemPhyXMini0424.q3_for_coupled_carnot_engine_and_refrigerator}
\uses{def:physics:phyx-mini-0424:phyxminiproblems-problemphyxmini0424-coupledcarnotsetup, def:physics:phyx-mini-0424:phyxminiproblems-problemphyxmini0424-heatinjoules, def:physics:phyx-mini-0424:phyxminiproblems-problemphyxmini0424-matchescoupledcarnotscenario, def:physics:phyx-mini-0424:phyxminiproblems-problemphyxmini0424-matchessuppliedfigure, def:physics:phyx-mini-0424:phyxminiproblems-problemphyxmini0424-hasphysicalmagnitudes, def:physics:phyx-mini-0424:phyxminiproblems-problemphyxmini0424-satisfiescoupledcarnotlaws}
The \(600/300\,\mathrm K\) Carnot engine drives the
\(500/400\,\mathrm K\) Carnot refrigerator losslessly.  With
\(Q_1=1000\,\mathrm J\), the rejected refrigerator heat is
\(Q_3=2500\,\mathrm J\), answer D.
\end{theorem}
\begin{proof}
Reversibility of the engine gives
\(300Q_1=600Q_2\), hence \(Q_2=500\,\mathrm J\).  Its first law therefore
gives \(W_{\rm out}=500\,\mathrm J\), and lossless coupling makes
\(W_{\rm in}=500\,\mathrm J\).  For the refrigerator,
\(400Q_3=500Q_4\) and \(Q_3=Q_4+500\).  Eliminating \(Q_4\) yields
\(Q_3=2500\,\mathrm J\).
\end{proof}

% --- Archon declaration topology begin ---
\section{Declaration topology}
\begin{definition}[energy Readout]
  \label{def:physics:phyx-mini-0424:phyxminiproblems-problemphyxmini0424-energyreadout}
  \lean{PhyXMiniProblems.ProblemPhyXMini0424.energyReadout}
  Read a dimensionful energy in coherent SI energy units
\end{definition}
\begin{proof}
  This is the declared model definition of energy Readout.
\end{proof}
\begin{definition}[energy In Joules]
  \label{def:physics:phyx-mini-0424:phyxminiproblems-problemphyxmini0424-energyinjoules}
  \lean{PhyXMiniProblems.ProblemPhyXMini0424.energyInJoules}
  \uses{def:physics:phyx-mini-0424:phyxminiproblems-problemphyxmini0424-energyreadout}
  Read heat or work in joules, normalized by the dimensionful one joule
\end{definition}
\begin{proof}
  This is the declared model definition of energy In Joules.
\end{proof}
\begin{definition}[temperature In Kelvins]
  \label{def:physics:phyx-mini-0424:phyxminiproblems-problemphyxmini0424-temperatureinkelvins}
  \lean{PhyXMiniProblems.ProblemPhyXMini0424.temperatureInKelvins}
  Read an absolute temperature in kelvin, accounting for the zero-preserving temperature unit in which the Physlib Temperature is stored
\end{definition}
\begin{proof}
  This is the declared model definition of temperature In Kelvins.
\end{proof}
\begin{definition}[Device]
  \label{def:physics:phyx-mini-0424:phyxminiproblems-problemphyxmini0424-device}
  \lean{PhyXMiniProblems.ProblemPhyXMini0424.Device}
  The two thermodynamic devices connected by the horizontal work shaft
\end{definition}
\begin{proof}
  This is the declared model definition of Device.
\end{proof}
\begin{definition}[Device Mode]
  \label{def:physics:phyx-mini-0424:phyxminiproblems-problemphyxmini0424-devicemode}
  \lean{PhyXMiniProblems.ProblemPhyXMini0424.DeviceMode}
  The operating role written next to each device in the figure
\end{definition}
\begin{proof}
  This is the declared model definition of Device Mode.
\end{proof}
\begin{definition}[Cycle Idealization]
  \label{def:physics:phyx-mini-0424:phyxminiproblems-problemphyxmini0424-cycleidealization}
  \lean{PhyXMiniProblems.ProblemPhyXMini0424.CycleIdealization}
  Whether a device is the reversible Carnot idealization stated in the problem
\end{definition}
\begin{proof}
  This is the declared model definition of Cycle Idealization.
\end{proof}
\begin{definition}[Device Position]
  \label{def:physics:phyx-mini-0424:phyxminiproblems-problemphyxmini0424-deviceposition}
  \lean{PhyXMiniProblems.ProblemPhyXMini0424.DevicePosition}
  The left and right device positions in the supplied diagram
\end{definition}
\begin{proof}
  This is the declared model definition of Device Position.
\end{proof}
\begin{definition}[Thermal Reservoir]
  \label{def:physics:phyx-mini-0424:phyxminiproblems-problemphyxmini0424-thermalreservoir}
  \lean{PhyXMiniProblems.ProblemPhyXMini0424.ThermalReservoir}
  The four reservoirs attached to the two devices
\end{definition}
\begin{proof}
  This is the declared model definition of Thermal Reservoir.
\end{proof}
\begin{definition}[Reservoir Position]
  \label{def:physics:phyx-mini-0424:phyxminiproblems-problemphyxmini0424-reservoirposition}
  \lean{PhyXMiniProblems.ProblemPhyXMini0424.ReservoirPosition}
  The four geometric reservoir positions visible in the diagram
\end{definition}
\begin{proof}
  This is the declared model definition of Reservoir Position.
\end{proof}
\begin{definition}[Heat Label]
  \label{def:physics:phyx-mini-0424:phyxminiproblems-problemphyxmini0424-heatlabel}
  \lean{PhyXMiniProblems.ProblemPhyXMini0424.HeatLabel}
  Heat-transfer symbols printed next to the four vertical arrows
\end{definition}
\begin{proof}
  This is the declared model definition of Heat Label.
\end{proof}
\begin{definition}[Work Label]
  \label{def:physics:phyx-mini-0424:phyxminiproblems-problemphyxmini0424-worklabel}
  \lean{PhyXMiniProblems.ProblemPhyXMini0424.WorkLabel}
  Work-transfer symbols printed on the horizontal connection
\end{definition}
\begin{proof}
  This is the declared model definition of Work Label.
\end{proof}
\begin{definition}[Heat Arrow Role]
  \label{def:physics:phyx-mini-0424:phyxminiproblems-problemphyxmini0424-heatarrowrole}
  \lean{PhyXMiniProblems.ProblemPhyXMini0424.HeatArrowRole}
  Physical direction and role of each heat arrow in the primary raster
\end{definition}
\begin{proof}
  This is the declared model definition of Heat Arrow Role.
\end{proof}
\begin{definition}[expected Heat Arrow Role]
  \label{def:physics:phyx-mini-0424:phyxminiproblems-problemphyxmini0424-expectedheatarrowrole}
  \lean{PhyXMiniProblems.ProblemPhyXMini0424.expectedHeatArrowRole}
  \uses{def:physics:phyx-mini-0424:phyxminiproblems-problemphyxmini0424-heatlabel, def:physics:phyx-mini-0424:phyxminiproblems-problemphyxmini0424-heatarrowrole}
  The heat-arrow roles associated with the four printed symbols
\end{definition}
\begin{proof}
  This is the declared model definition of expected Heat Arrow Role.
\end{proof}
\begin{definition}[Work Arrow Role]
  \label{def:physics:phyx-mini-0424:phyxminiproblems-problemphyxmini0424-workarrowrole}
  \lean{PhyXMiniProblems.ProblemPhyXMini0424.WorkArrowRole}
  Physical role of each work label on the shared horizontal shaft
\end{definition}
\begin{proof}
  This is the declared model definition of Work Arrow Role.
\end{proof}
\begin{definition}[expected Work Arrow Role]
  \label{def:physics:phyx-mini-0424:phyxminiproblems-problemphyxmini0424-expectedworkarrowrole}
  \lean{PhyXMiniProblems.ProblemPhyXMini0424.expectedWorkArrowRole}
  \uses{def:physics:phyx-mini-0424:phyxminiproblems-problemphyxmini0424-worklabel, def:physics:phyx-mini-0424:phyxminiproblems-problemphyxmini0424-workarrowrole}
  The work-arrow roles associated with the two printed symbols
\end{definition}
\begin{proof}
  This is the declared model definition of expected Work Arrow Role.
\end{proof}
\begin{definition}[Coupled Carnot Figure]
  \label{def:physics:phyx-mini-0424:phyxminiproblems-problemphyxmini0424-coupledcarnotfigure}
  \lean{PhyXMiniProblems.ProblemPhyXMini0424.CoupledCarnotFigure}
  \uses{def:physics:phyx-mini-0424:phyxminiproblems-problemphyxmini0424-device, def:physics:phyx-mini-0424:phyxminiproblems-problemphyxmini0424-devicemode, def:physics:phyx-mini-0424:phyxminiproblems-problemphyxmini0424-deviceposition, def:physics:phyx-mini-0424:phyxminiproblems-problemphyxmini0424-thermalreservoir, def:physics:phyx-mini-0424:phyxminiproblems-problemphyxmini0424-reservoirposition, def:physics:phyx-mini-0424:phyxminiproblems-problemphyxmini0424-heatlabel, def:physics:phyx-mini-0424:phyxminiproblems-problemphyxmini0424-worklabel, def:physics:phyx-mini-0424:phyxminiproblems-problemphyxmini0424-heatarrowrole, def:physics:phyx-mini-0424:phyxminiproblems-problemphyxmini0424-workarrowrole}
  Qualitative information displayed by the supplied raster. Numerical temperatures and energy readouts are related to the setup separately in MatchesSuppliedFigure
\end{definition}
\begin{proof}
  This is the declared model definition of Coupled Carnot Figure.
\end{proof}
\begin{definition}[Coupled Carnot Setup]
  \label{def:physics:phyx-mini-0424:phyxminiproblems-problemphyxmini0424-coupledcarnotsetup}
  \lean{PhyXMiniProblems.ProblemPhyXMini0424.CoupledCarnotSetup}
  \uses{def:physics:phyx-mini-0424:phyxminiproblems-problemphyxmini0424-device, def:physics:phyx-mini-0424:phyxminiproblems-problemphyxmini0424-devicemode, def:physics:phyx-mini-0424:phyxminiproblems-problemphyxmini0424-cycleidealization, def:physics:phyx-mini-0424:phyxminiproblems-problemphyxmini0424-thermalreservoir, def:physics:phyx-mini-0424:phyxminiproblems-problemphyxmini0424-heatlabel, def:physics:phyx-mini-0424:phyxminiproblems-problemphyxmini0424-worklabel, def:physics:phyx-mini-0424:phyxminiproblems-problemphyxmini0424-coupledcarnotfigure}
  The physical observables of the two devices. Heat and work fields are independent dimensionful magnitudes; none is defined from the requested answer
\end{definition}
\begin{proof}
  This is the declared model definition of Coupled Carnot Setup.
\end{proof}
\begin{definition}[heat In Joules]
  \label{def:physics:phyx-mini-0424:phyxminiproblems-problemphyxmini0424-heatinjoules}
  \lean{PhyXMiniProblems.ProblemPhyXMini0424.heatInJoules}
  \uses{def:physics:phyx-mini-0424:phyxminiproblems-problemphyxmini0424-energyinjoules, def:physics:phyx-mini-0424:phyxminiproblems-problemphyxmini0424-heatlabel, def:physics:phyx-mini-0424:phyxminiproblems-problemphyxmini0424-coupledcarnotsetup}
  Joule readout of one of the four labelled heat-transfer magnitudes
\end{definition}
\begin{proof}
  This is the declared model definition of heat In Joules.
\end{proof}
\begin{definition}[work In Joules]
  \label{def:physics:phyx-mini-0424:phyxminiproblems-problemphyxmini0424-workinjoules}
  \lean{PhyXMiniProblems.ProblemPhyXMini0424.workInJoules}
  \uses{def:physics:phyx-mini-0424:phyxminiproblems-problemphyxmini0424-energyinjoules, def:physics:phyx-mini-0424:phyxminiproblems-problemphyxmini0424-worklabel, def:physics:phyx-mini-0424:phyxminiproblems-problemphyxmini0424-coupledcarnotsetup}
  Joule readout of one of the two labelled work-transfer magnitudes
\end{definition}
\begin{proof}
  This is the declared model definition of work In Joules.
\end{proof}
\begin{definition}[reservoir Temperature In Kelvins]
  \label{def:physics:phyx-mini-0424:phyxminiproblems-problemphyxmini0424-reservoirtemperatureinkelvins}
  \lean{PhyXMiniProblems.ProblemPhyXMini0424.reservoirTemperatureInKelvins}
  \uses{def:physics:phyx-mini-0424:phyxminiproblems-problemphyxmini0424-temperatureinkelvins, def:physics:phyx-mini-0424:phyxminiproblems-problemphyxmini0424-thermalreservoir, def:physics:phyx-mini-0424:phyxminiproblems-problemphyxmini0424-coupledcarnotsetup}
  Kelvin readout of a labelled reservoir's absolute temperature
\end{definition}
\begin{proof}
  This is the declared model definition of reservoir Temperature In Kelvins.
\end{proof}
\begin{definition}[Matches Coupled Carnot Scenario]
  \label{def:physics:phyx-mini-0424:phyxminiproblems-problemphyxmini0424-matchescoupledcarnotscenario}
  \lean{PhyXMiniProblems.ProblemPhyXMini0424.MatchesCoupledCarnotScenario}
  \uses{def:physics:phyx-mini-0424:phyxminiproblems-problemphyxmini0424-coupledcarnotsetup}
  The prose-level scenario: a Carnot heat engine directly drives a Carnot refrigerator. Quantitative energy and reversibility laws are stated separately
\end{definition}
\begin{proof}
  This is the declared model definition of Matches Coupled Carnot Scenario.
\end{proof}
\begin{definition}[Matches Supplied Figure]
  \label{def:physics:phyx-mini-0424:phyxminiproblems-problemphyxmini0424-matchessuppliedfigure}
  \lean{PhyXMiniProblems.ProblemPhyXMini0424.MatchesSuppliedFigure}
  \uses{def:physics:phyx-mini-0424:phyxminiproblems-problemphyxmini0424-expectedheatarrowrole, def:physics:phyx-mini-0424:phyxminiproblems-problemphyxmini0424-expectedworkarrowrole, def:physics:phyx-mini-0424:phyxminiproblems-problemphyxmini0424-coupledcarnotsetup, def:physics:phyx-mini-0424:phyxminiproblems-problemphyxmini0424-heatinjoules, def:physics:phyx-mini-0424:phyxminiproblems-problemphyxmini0424-reservoirtemperatureinkelvins}
  Facts read from the primary image. They include positions, arrows, the four reservoir temperatures, and Q₁ = 1000 J, but no numerical value for Q₃, Q₄, or either work transfer
\end{definition}
\begin{proof}
  This is the declared model definition of Matches Supplied Figure.
\end{proof}
\begin{definition}[Has Physical Magnitudes]
  \label{def:physics:phyx-mini-0424:phyxminiproblems-problemphyxmini0424-hasphysicalmagnitudes}
  \lean{PhyXMiniProblems.ProblemPhyXMini0424.HasPhysicalMagnitudes}
  \uses{def:physics:phyx-mini-0424:phyxminiproblems-problemphyxmini0424-coupledcarnotsetup, def:physics:phyx-mini-0424:phyxminiproblems-problemphyxmini0424-heatinjoules, def:physics:phyx-mini-0424:phyxminiproblems-problemphyxmini0424-workinjoules, def:physics:phyx-mini-0424:phyxminiproblems-problemphyxmini0424-reservoirtemperatureinkelvins}
  Positivity conditions for the directed heat and work magnitudes
\end{definition}
\begin{proof}
  This is the declared model definition of Has Physical Magnitudes.
\end{proof}
\begin{definition}[Satisfies Coupled Carnot Laws]
  \label{def:physics:phyx-mini-0424:phyxminiproblems-problemphyxmini0424-satisfiescoupledcarnotlaws}
  \lean{PhyXMiniProblems.ProblemPhyXMini0424.SatisfiesCoupledCarnotLaws}
  \uses{def:physics:phyx-mini-0424:phyxminiproblems-problemphyxmini0424-coupledcarnotsetup, def:physics:phyx-mini-0424:phyxminiproblems-problemphyxmini0424-heatinjoules, def:physics:phyx-mini-0424:phyxminiproblems-problemphyxmini0424-workinjoules, def:physics:phyx-mini-0424:phyxminiproblems-problemphyxmini0424-reservoirtemperatureinkelvins}
  Governing relations for the ideal coupled devices: the engine obeys Q₁ = Q₂ + W\_out; reversibility gives Q₁ / T\_engine,hot = Q₂ / T\_engine,cold; the refrigerator obeys Q₃ = Q₄ + W\_in; reversibility gives Q₃ / T\_refrigerator,hot = Q₄ / T\_refrigerator,cold; and the lossless direct shaft gives W\_out = W\_in. The reversible heat ratios are written without division as cross-products of kelvin and joule readouts. None of these general laws contains the requested numerical value of Q₃
\end{definition}
\begin{proof}
  This is the declared model definition of Satisfies Coupled Carnot Laws.
\end{proof}
\begin{definition}[Answer Choice]
  \label{def:physics:phyx-mini-0424:phyxminiproblems-problemphyxmini0424-answerchoice}
  \lean{PhyXMiniProblems.ProblemPhyXMini0424.AnswerChoice}
  The four answer labels supplied with the problem
\end{definition}
\begin{proof}
  This is the declared model definition of Answer Choice.
\end{proof}
\begin{definition}[displayed Q3 In Joules]
  \label{def:physics:phyx-mini-0424:phyxminiproblems-problemphyxmini0424-displayedq3injoules}
  \lean{PhyXMiniProblems.ProblemPhyXMini0424.displayedQ3InJoules}
  \uses{def:physics:phyx-mini-0424:phyxminiproblems-problemphyxmini0424-answerchoice}
  The candidate Q₃ magnitudes printed beside the answer labels, in joules
\end{definition}
\begin{proof}
  This is the declared model definition of displayed Q3 In Joules.
\end{proof}
\begin{definition}[recorded Dataset Answer]
  \label{def:physics:phyx-mini-0424:phyxminiproblems-problemphyxmini0424-recordeddatasetanswer}
  \lean{PhyXMiniProblems.ProblemPhyXMini0424.recordedDatasetAnswer}
  \uses{def:physics:phyx-mini-0424:phyxminiproblems-problemphyxmini0424-answerchoice}
  The source dataset records answer label D; this metadata is not a premise
\end{definition}
\begin{proof}
  This is the declared model definition of recorded Dataset Answer.
\end{proof}
% --- Archon declaration topology end ---
% --- Archon physics formalization source end ---
